% ==========================================
% LEXIQUE
% ==========================================
\chapter*{Lexique}
\addcontentsline{toc}{chapter}{Lexique}

Ce lexique regroupe les termes métier, acronymes internes et concepts spécifiques au projet utilisés dans ce rapport.

\vspace{0.5cm}

\section*{Termes métier assurance}

\begin{description}
    \item[AFN (Affaires Nouvelles)] : Polices souscrites ou renouvelées au cours de la période. Indicateur de production commerciale.
    \item[AZ] : Canal de distribution regroupant les Agents et le Courtage traditionnel. Se distingue d'AZEC par ses sources de données et sa logique métier.
    \item[AZEC] : Canal de Courtage Construction (entreprises de construction). Sources de données et règles de gestion distinctes du canal AZ.
    \item[Coassurance] : Partage d'un risque entre plusieurs assureurs. La compagnie ne porte qu'une fraction du risque total (en pourcentage).
    \item[LCI (Limite Contractuelle d'Indemnité)] : Montant maximum que l'assureur s'engage à verser en cas de sinistre.
    \item[PE (Perte d'Exploitation)] : Capital assuré couvrant les pertes financières consécutives à un sinistre.
    \item[PTF (Portefeuille)] : Ensemble des contrats ou polices d'assurance actifs, regroupant les informations sur les clients, les couvertures, et les risques assurés à un moment donné.
    \item[RD (Risque Direct)] : Capital assuré pour les dommages matériels directs.
    \item[RES (Résiliations)] : Polices résiliées ou non renouvelées au cours de la période.
    \item[RPC (Remises en Portefeuille Avec Modification de Prime)] : Polices réintégrées avec ajustement tarifaire.
    \item[RPT (Remises en Portefeuille Sans Modification de Prime)] : Polices réintégrées sans changement de prime.
    \item[SMP (Sinistre Maximum Possible)] : Estimation du coût maximum d'un sinistre pour l'assureur.
\end{description}

\section*{Acronymes internes Allianz}

\begin{description}
    \item[ADC (Advanced Data \& Climate)] : Pôle d'Allianz France spécialisé en data management et modélisations. Équipe d'accueil du projet.
    \item[CDPOLE] : Code Pôle de distribution : « 1 » = Agents, « 3 » = Courtage.
    \item[DIRCOM] : Direction commerciale / canal de distribution (« AZ » ou « AZEC »).
    \item[IRD (Inventaire des Risques et Données)] : Référentiel Allianz décrivant les caractéristiques techniques des risques assurés.
    \item[P4D] : Direction d'Allianz France regroupant plusieurs pôles data, dont ADC.
    \item[SAS EXIST] : Programme stratégique de migration des datamarts SAS vers Azure/PySpark.
\end{description}


